\documentclass[10pt]{article}

% ---------- Document Identity (update these only) ----------
\newcommand{\DocStage}{}
\newcommand{\DocTitle}{Your Road to Tier One SOF}
\newcommand{\ProgramName}{182-Week Civilian-to-Selection Readiness Pipeline}
\newcommand{\ProgramSubtitle}{}

% ---------- Page ----------
\usepackage[legalpaper,left=0.5in,right=0.5in,top=0.6in,bottom=0.45in]{geometry}

% ---------- Typography ----------
\usepackage[T1]{fontenc}
\usepackage[utf8]{inputenc}
\usepackage{libertinus}
\usepackage{microtype}
\usepackage{parskip}
\usepackage{ragged2e}
\usepackage{textcomp}
\AtBeginDocument{\color{black}}
\AtBeginDocument{\pagecolor{white}}
\AtBeginDocument{\justifying}

% ---------- Landscape pages ----------
\usepackage{pdflscape}

% ---------- Color ----------
\usepackage[table]{xcolor}
\definecolor{StageBlue}{HTML}{1F4E79}
\definecolor{StageBlueDark}{HTML}{163A5A}
\definecolor{LightGray}{HTML}{F2F4F7}
\definecolor{MidGray}{HTML}{D9DEE5}
\definecolor{RuleGray}{HTML}{C7CED8}
\definecolor{HeaderInk}{HTML}{000000}
\definecolor{Ink}{HTML}{000000}

% ---------- Utility ----------
\usepackage{enumitem}
\sloppy
\setlength{\emergencystretch}{2em}

% ---------- Boxes / UI ----------
\usepackage[most]{tcolorbox}

\tcbset{
  charterTitle/.style={
    enhanced,
    colback=StageBlue!4,
    colframe=StageBlue,
    boxrule=1.25pt,
    arc=3pt,
    left=16pt,right=16pt,top=14pt,bottom=14pt,
    borderline north={3.2pt}{0pt}{StageBlue},
    borderline west={4pt}{0pt}{StageBlue},
    borderline south={1.25pt}{0pt}{StageBlue},
    drop shadow={black!25!white},
  },
  charterRule/.style={
    enhanced,
    colback=LightGray,
    colframe=RuleGray,
    boxrule=0.85pt,
    arc=3pt,
    left=12pt,right=12pt,top=10pt,bottom=10pt,
    borderline west={3pt}{0pt}{StageBlue},
  },
  charterBadge/.style={
    enhanced,
    colback=StageBlue,
    coltext=white,
    colframe=StageBlueDark,
    boxrule=0.6pt,
    arc=2pt,
    left=6pt,right=6pt,top=3pt,bottom=3pt,
  }
}
\newtcbox{\Badge}{nobeforeafter,charterBadge}

% ---------- Lists ----------
\setlist[itemize]{leftmargin=1.5em, itemsep=0.25em, topsep=0.25em}
\setlist[enumerate]{leftmargin=1.8em, itemsep=0.25em, topsep=0.25em}

% ---------- TOC ----------
\usepackage{tocloft}
\setcounter{tocdepth}{3}
\setlength{\cftbeforesecskip}{3pt}
\setlength{\cftbeforesubsecskip}{2pt}
\setlength{\cftbeforesubsubsecskip}{0pt}
\renewcommand{\cftsecfont}{\bfseries}
\renewcommand{\cftsecpagefont}{\bfseries}
\renewcommand{\cftsubsecfont}{\normalfont}
\renewcommand{\cftsubsecpagefont}{\normalfont}
\renewcommand{\cftsubsubsecfont}{\normalfont}
\renewcommand{\cftsubsubsecpagefont}{\normalfont}
\setlength{\cftsecindent}{0pt}
\setlength{\cftsubsecindent}{1.25em}
\setlength{\cftsubsubsecindent}{2.8em}
\setlength{\cftsecnumwidth}{2.1em}
\setlength{\cftsubsecnumwidth}{2.8em}
\setlength{\cftsubsubsecnumwidth}{3.6em}

% Remove the built-in TOC heading so only the custom heading is shown
\makeatletter
\renewcommand{\tableofcontents}{\@starttoc{toc}}
\makeatother

% ---------- Header/Footer: page number only ----------
\usepackage{fancyhdr}
\pagestyle{fancy}
\fancyhf{}
\renewcommand{\headrulewidth}{0pt}
\renewcommand{\footrulewidth}{0pt}
\fancyfoot[C]{\thepage}

% ---------- Section formatting ----------
\usepackage{titlesec}

\titleformat{\section}
  {\Large\bfseries\color{Ink}}
  {\thesection}{0.6em}{}
  [\vspace{0.25em}\textcolor{StageBlue}{\rule{\linewidth}{0.8pt}}]

\titleformat{\subsection}
  {\large\bfseries\color{Ink}}
  {\thesubsection}{0.6em}{}

\titleformat{\subsubsection}
  {\normalsize\bfseries\color{Ink}}
  {\thesubsubsection}{0.6em}{}

\titlespacing*{\section}{0pt}{0.95em}{0.35em}
\titlespacing*{\subsection}{0pt}{0.65em}{0.25em}
\titlespacing*{\subsubsection}{0pt}{0.55em}{0.2em}

% ---------- Table engine ----------
\usepackage{tabularray}
\UseTblrLibrary{booktabs}

% ---------- tabularray: GLOBAL table treatment ----------
\SetTblrInner{
  cells = {fg=black, bg=white, valign=t},
}

% ---------- Header helpers ----------
\newcommand{\Hdr}[1]{\shortstack[c]{\strut #1\strut}}

% ---------- Lines ----------
\newcommand{\TitleLine}{\textcolor{StageBlue}{\rule{0.98\linewidth}{0.95pt}}}
\newcommand{\SubTitleLine}{\textcolor{RuleGray}{\rule{0.98\linewidth}{0.65pt}}}

% ---------- Continued on next page ----------
\newcommand{\ContinuedOnNextPageInline}{%
  \par
  \vfill
  \noindent\textcolor{RuleGray}{\rule{\linewidth}{1pt}}\vspace{-2pt}
  \noindent\hfill\textit{Continued on next page}\par\vspace{-10pt}
  \noindent\textcolor{RuleGray}{\rule{\linewidth}{1pt}}
  \clearpage
}

% ---------- PDF hyperlinks ----------
\usepackage[hidelinks]{hyperref}
\hypersetup{
  colorlinks=false,
  pdfborder={0 0 0},
  linktoc=all,
  pdftitle={\DocTitle},
  pdfauthor={\ProgramName}
}

% =========================================================
\begin{document}

\subsection*{Phase 08 Card Header}
\phantomsection
\addcontentsline{toc}{subsection}{Phase 08 Card Header}

\begin{itemize}
\item Phase \textbf{ID}: Phase 08
\item Phase Name: Structural Durability and Load Tolerance
\item Duration weeks: 12
\item Version: v1.0
\item Stage Alignment: Stage 5
\item Governing Status: Deployable
\end{itemize}

\subsection*{1. Phase Mission}
\phantomsection


Build load-bearing and interface durability (feet, skin, gait mechanics, pack interface) so long-duration ruck milestones can be executed safely and produce admissible evidence without compromising \textbf{RUN} and \textbf{CAL} validity.

\subsection*{2. Phase Purpose}
\phantomsection


\begin{itemize}
\item Establish resilient tissue tolerance and equipment-interface stability under high time-on-feet, with strict controls to prevent blisters, gait drift, and overuse cascades.
\item Set up Phase 09 by proving the trainee can maintain durability and recovery stability while handling high exposure density without invalidating evidence.
\item Reinforce anti-stacking discipline: Phase 08 treats the long ruck milestone as the dominant high-cost proof event and stages other tests around it.
\end{itemize}

\subsection*{3. Entry Criteria}
\phantomsection


\begin{itemize}
\item Prior phase exit posture is satisfied with admissible evidence; default posture is \textbf{HOLD} if evidence is incomplete or invalid.
\item Decision-window compliance posture is met under Stage 1.F (no reliance on inadmissible or incomplete logs).
\item Foot/skin interface is stable enough to attempt load-tolerance work:
  \begin{itemize}
  \item no active blister,
  \item no unresolved gait-altering pain,
  \item hotspot protocol executed and trending stable.
  \end{itemize}
\item Required Phase 08 category capabilities are in place before attempting long ruck milestones:
  \begin{itemize}
  \item 7.07 and 7.11 readiness for long ruck exposure,
  \item 7.19 protective/assistive readiness for interface resilience,
  \item 7.18 timing/measurement readiness for admissible test evidence.
  \end{itemize}
\end{itemize}

\subsection*{4. Operating Constraints}
\phantomsection


\begin{itemize}
\item 6 sessions per week.
\item Required session elements per Stage 1.F in correct order:
  \begin{itemize}
  \item Safety self-check first
  \item Warm-up
  \item Mobility
  \item Main work
  \item Cardio
  \item Cooldown
  \item Recovery notes and any required post-session notes per Stage 1.F
  \end{itemize}
\item Cardio minimum standard: minimum 30 minutes continuous per session.
\item 10,000 steps per day global requirement.
\item Validity doctrine: admissible Valid evidence only for gates; \textbf{INVALID} provides no gate credit and triggers reslot posture.
\item Water governance remains absolute: no unsupervised underwater or breath-hold exposure ever; underwater is not a Phase 08 gate requirement.
\end{itemize}

\subsection*{5. Primary Adaptations and Physiological Focus}
\phantomsection


\begin{itemize}
\item Load-tolerance and interface hardening: feet, skin, shoulders/traps, and gait mechanics under sustained carriage.
\item Durability stability: reduce hotspots/blisters and prevent gait drift under long exposure.
\item Recovery resilience: maintain sleep stability and soreness/fatigue control while exposure cost rises.
\item Validity stability: preserve comparable routes/tools and protect test windows from novelty and fatigue contamination.
\end{itemize}

\subsection*{6. Primary Modalities}
\phantomsection


\begin{itemize}
\item Emphasized domains: \textbf{RUCK}, \textbf{DURABILITY}, \textbf{RECOVERY}, \textbf{RUN}, \textbf{CAL}, \textbf{BODYCOMP}.
\item Supporting domains: \textbf{SWIM} (surface-only unit-locked test), \textbf{STR} (maintenance only; no new risk escalation), \textbf{AER} (only as low-impact substitution when needed; not interchangeable with \textbf{RUN} evidence).
\item Prohibited or de-emphasized:
  \begin{itemize}
  \item no new underwater gate objectives (governance-gated and not required here),
  \item no impact escalation by “testing harder” outside protected windows,
  \item no ruck load escalation outside the phase’s controlled posture and evidence windows.
  \end{itemize}
\end{itemize}

\clearpage
\begin{landscape}

\subsection*{7. Weekly Structure}
\phantomsection


\begingroup
\scriptsize
\setlength{\tabcolsep}{2pt}
\renewcommand{\arraystretch}{1.12}

\begin{longtblr}{
  width=\linewidth,
  rowhead=1,
  colspec = {
    Q[l,wd=0.70in]
    X[l,1.60]
    X[l,1.15]
    X[l,2.80]
    X[l,2.70]
  },
  hlines,
  vlines,
  cells = {hsep=6pt, vsep=6pt, halign=l, valign=t},
  row{odd} = {bg=white},
  row{even} = {bg=LightGray},
  row{1} = {bg=StageBlue!15, fg=black, font=\bfseries, halign=l, valign=m},
}
\Hdr{Session \#} &
\Hdr{Primary Focus} &
\Hdr{Main Work Domains} &
\Hdr{Cardio Modality/Intent} &
\Hdr{Notes/Risk Controls} \\

1 &
Durability interface hardening &
\textbf{RUCK}; \textbf{DURABILITY} &
Modality: step-based or treadmill walk; Minutes: 30–60; RPE plus Talk Test: RPE 3–5, full-to-short sentences; Continuity: continuous planned &
Feet/skin check pre/post; no novelty in footwear; hotspots trigger downshift next session and documentation; avoid high-risk terrain. \\

2 &
\textbf{CAL} repeatability + trunk control &
\textbf{CAL}; \textbf{DURABILITY} &
Modality: low-impact steady; Minutes: 30–45; RPE plus Talk Test: RPE 2–4, full sentences; Continuity: continuous planned &
\textbf{CAL} is technique-governed; protect joints; avoid fatigue spillover into ruck days. \\

3 &
\textbf{RUN} durability check posture &
\textbf{RUN}; \textbf{DURABILITY} &
Modality: step-based or treadmill; Minutes: 30–50; RPE plus Talk Test: RPE 3–5, full-to-short sentences; Continuity: continuous planned &
No maximal \textbf{RUN} near long ruck; route/tool stability enforced; pain/mechanics drift triggers regression and reslot of any planned test. \\

4 &
Recovery-biased throughput &
\textbf{RECOVERY}; \textbf{DURABILITY} &
Modality: lowest-risk continuous; Minutes: 30–45; RPE plus Talk Test: RPE 2–3, full sentences; Continuity: continuous planned &
Minimum Viable Session candidate day; preserve compliance without adding fatigue cost. \\

5 &
Load-tolerance reinforcement &
\textbf{RUCK}; \textbf{DURABILITY} &
Modality: step-based; Minutes: 30–60; RPE plus Talk Test: RPE 3–5, full-to-short sentences; Continuity: continuous planned &
Treat as controlled exposure day; do not “make up” missed work by increasing intensity; protect feet and gait. \\

6 &
Swim and mobility integration &
\textbf{SWIM}; \textbf{RECOVERY} &
Modality: low-impact steady; Minutes: 30–45; RPE plus Talk Test: RPE 2–4, full sentences; Continuity: continuous planned &
Surface-only; pool verification required when testing; no underwater; reslot swim test if pool conditions compromise validity. \\

\end{longtblr}

\endgroup

\subsection*{8. Progression Rules}
\phantomsection


\begin{itemize}
\item One progression lever at a time:
  \begin{itemize}
  \item do not increase ruck load/tolerance, \textbf{RUN} stress, and \textbf{CAL} density in the same week.
  \end{itemize}
\item Long ruck is the dominant proof exposure:
  \begin{itemize}
  \item protect it by reducing competing high-cost demands in the lead-in and inside the same protected window.
  \end{itemize}
\item Regression triggers (immediate):
  \begin{itemize}
  \item pain alters mechanics, hotspot escalation, blister formation, unusual breathlessness, dizziness/lightheadedness, swelling/instability → downshift modality/intensity and document; reslot affected tests.
  \end{itemize}
\item Substitution posture:
  \begin{itemize}
  \item preserve intent and reduce risk; use \textbf{AER} modality substitution only to preserve continuous Cardio safely; do not claim \textbf{RUN} tier attainment from \textbf{AER} evidence.
  \end{itemize}
\item No stacking make-ups:
  \begin{itemize}
  \item missed or invalid tests are reslotted per Stage 3.E, not compressed.
  \end{itemize}
\end{itemize}

\subsection*{9. Deload and Test Placement}
\phantomsection


\begin{itemize}
\item Phase 08 taper-aware protected windows (week-in-phase):
  \begin{itemize}
  \item Week 1 baseline (low-risk touchpoint),
  \item Week 5 Deload+Test (mid-phase decision window),
  \item Week 10 Deload+Test (second decision window; long ruck milestone by intent),
  \item Week 11 Transition (taper/stabilization; no new maximal tests),
  \item Week 12 Deload+Test (end-phase exit gate window).
  \end{itemize}
\item If Stage 3 integrated calendars specify a different placement, Stage 3 controls final placement.
\end{itemize}

\end{landscape}

\clearpage
\begin{landscape}

\subsection*{10. Test Battery}
\phantomsection


\begingroup
\scriptsize
\setlength{\tabcolsep}{2pt}
\renewcommand{\arraystretch}{1.12}

\begin{longtblr}{
  width=\linewidth,
  rowhead=1,
  colspec = {
    X[l,1.55]
    X[l,2.80]
    Q[l,wd=0.90in]
    X[l,3.05]
    Q[l,wd=1.25in]
  },
  hlines,
  vlines,
  cells = {hsep=6pt, vsep=6pt, halign=l, valign=t},
  row{odd} = {bg=white},
  row{even} = {bg=LightGray},
  row{1} = {bg=StageBlue!15, fg=black, font=\bfseries, halign=l, valign=m},
}
\Hdr{Test Timing} &
\Hdr{Test \textbf{ID}/Name} &
\Hdr{Domain} &
\Hdr{Validity Conditions} &
\Hdr{Pass Threshold Stage 2 Tier} \\

Week 5 Deload+Test (mid-phase) &
\textbf{C12} / \textbf{MET-2B-019} 6-Mile Ruck Time Trial — 35 lb dry &
\textbf{RUCK} &
Dry load only (water logged separately); measured route or treadmill identity logged; footwear and route class stable; no stop-rule violations; complete Stage 2.E fields &
\textbf{INTERMEDIATE} \\

Week 5 Deload+Test (mid-phase) &
\textbf{C10} / \textbf{MET-2B-016} 3-Mile Run Time Trial &
\textbf{RUN} &
Measured route or treadmill identity/settings logged; warm-up logged; no stop-rule violations; no novelty near window; complete Stage 2.E fields &
\textbf{INTERMEDIATE} \\

Week 5 Deload+Test (mid-phase) &
\textbf{C4} / \textbf{MET-2B-010} Push-Ups — 2-Minute Timed, Strict Standard &
\textbf{CAL} &
\textbf{CAL} artifact evidence required for \textbf{VALID} gate use; strict standard; timing verifiable; complete Stage 2.E fields &
\textbf{INTERMEDIATE} \\

Week 5 Deload+Test (mid-phase) &
\textbf{C5} / \textbf{MET-2B-011} Sit-Ups — 2-Minute Timed, Standard &
\textbf{CAL} &
\textbf{CAL} artifact evidence required for \textbf{VALID} gate use; strict timing; complete Stage 2.E fields &
\textbf{INTERMEDIATE} \\

Week 5 Deload+Test (mid-phase) &
\textbf{C6} / \textbf{MET-2B-012} Pull-Ups — Strict Dead-Hang &
\textbf{CAL} &
\textbf{CAL} artifact evidence required for \textbf{VALID} gate use; strict ROM; complete Stage 2.E fields &
\textbf{INTERMEDIATE} \\

Week 5 Deload+Test (mid-phase) &
\textbf{C7} / \textbf{MET-2B-013} Plank — Forearm &
\textbf{CAL} &
\textbf{CAL} artifact evidence required for \textbf{VALID} gate use; strict termination; complete Stage 2.E fields &
\textbf{INTERMEDIATE} \\

Week 10 Deload+Test (milestone) &
\textbf{C12} / \textbf{MET-2B-021} 12-Mile Ruck Time Trial — 55 lb dry &
\textbf{RUCK} &
Dry load only (water logged separately); route/tool/footwear lock enforced; stop posture and blister governance enforced; complete Stage 2.E fields; treat as dominant stand-alone major event by intent &
\textbf{INTERMEDIATE} \\

Week 12 Deload+Test (end-phase) &
\textbf{C13} / \textbf{MET-2B-022} Swim 500 yd Time — No Fins &
\textbf{SWIM} &
Pool unit verified; no fins; allowable strokes; warm-up logged; no standing on bottom or hanging on wall/line; complete Stage 2.E fields &
\textbf{INTERMEDIATE} \\

Week 12 Deload+Test (end-phase) &
\textbf{C12} / \textbf{MET-2B-019} 6-Mile Ruck Time Trial — 35 lb dry &
\textbf{RUCK} &
Same validity posture as Week 5; staged so it does not compromise swim validity; reslot if spacing cannot be protected &
\textbf{INTERMEDIATE} \\

Week 12 Deload+Test (end-phase) &
\textbf{C10} / \textbf{MET-2B-016} 3-Mile Run Time Trial &
\textbf{RUN} &
Only if spacing and durability stability permit; otherwise omit and reslot to next protected window; validity posture per Stage 2 and Stage 2.E &
\textbf{INTERMEDIATE} \\

Week 12 Deload+Test (end-phase) &
\textbf{C16} / \textbf{MET-2B-036} Gait-Altering Pain Flag (plus \textbf{MET-2B-037} foot flag) &
\textbf{DURABILITY} &
Must remain No for advancement posture; complete logs; any Yes constrains gate outcome &
No / Not Active \\

\end{longtblr}

\endgroup

\end{landscape}

\clearpage

\subsection*{11. Exit Standards}
\phantomsection


\begin{itemize}
\item Phase 08 exit posture is \textbf{INTERMEDIATE} at the end gate and is evaluated only using admissible evidence:
  \begin{itemize}
  \item \textbf{VALID} tests in the Week 12 Deload+Test window for the listed gate items, and
  \item decision-window compliance posture per Stage 1.F.
  \end{itemize}
\item Minimum exit conditions (objective gate posture):
  \begin{itemize}
  \item \textbf{RUCK}: \textbf{C12} / \textbf{MET-2B-021} executed as a \textbf{VALID} milestone within Phase 08 (Week 10 is the intended window). If the milestone is \textbf{INVALID} or not executed, default posture is \textbf{HOLD} and reslot to the next available protected window (or extend phase under change control if needed).
  \item \textbf{DURABILITY}: \textbf{MET-2B-036} gait-altering pain flag remains No for advancement posture; foot hotspot/blister posture stable and documented.
  \item \textbf{SWIM}: \textbf{C13} / \textbf{MET-2B-022} or \textbf{MET-2B-023} executed as \textbf{VALID} in the end gate window where phase contract requires it; if pool validity is compromised, reslot.
  \item \textbf{CAL}: battery items remain \textbf{VALID} where executed and do not degrade durability posture; invalid \textbf{CAL} evidence is excluded from gate claims.
  \end{itemize}
\item Explicit admissibility: any \textbf{INVALID} test provides no gate credit and cannot support \textbf{ADVANCE}; reslot per Stage 3.E.
\end{itemize}

\subsection*{12. Risk Controls and Stop Rules}
\phantomsection


\begin{itemize}
\item Authority references: Stage 0.5 boundaries and Stage 1.F stop/validity doctrine control; Stage 3.E controls reslot and phase extension posture.
\item Phase 08 high-risk exposure controls:
  \begin{itemize}
  \item Long ruck is the dominant high-cost exposure; it must be treated as stand-alone by intent and protected from adjacent maximal \textbf{RUN}.
  \item Foot/skin governance is decisive: hotspots trending to blisters trigger downshift; active blister blocks long ruck gate attempts until resolved.
  \item Environment hazard controls (ice, heat, wind, low visibility) trigger substitution or reslot; compromised attempts are \textbf{INVALID} for gate use.
  \item No breathless intensity chasing; unusual/disproportionate breathlessness triggers immediate downshift or \textbf{STOP} \textbf{NOW} posture per Stage 1.F.
  \end{itemize}
\end{itemize}

\clearpage
\begin{landscape}

\subsection*{13. Required Capabilities and Dependencies}
\phantomsection


\begingroup
\scriptsize
\setlength{\tabcolsep}{2pt}
\renewcommand{\arraystretch}{1.12}

\begin{longtblr}{
  width=\linewidth,
  rowhead=1,
  colspec = {
    Q[l,wd=0.85in]
    X[l,2.95]
    Q[l,wd=1.35in]
    Q[l,wd=1.20in]
    X[l,3.05]
  },
  hlines,
  vlines,
  cells = {hsep=6pt, vsep=6pt, halign=l, valign=t},
  row{odd} = {bg=white},
  row{even} = {bg=LightGray},
  row{1} = {bg=StageBlue!15, fg=black, font=\bfseries, halign=l, valign=m},
}
\Hdr{Dependency \textbf{ID}} &
\Hdr{Required Capability Category and Tier} &
\Hdr{Due-By week-in-phase} &
\Hdr{Owner} &
\Hdr{Fail Action} \\

\textbf{DEP-1C-006} &
7.18 Testing Timing Measurement and Course-Setup Tools — Tier 3 (lock-in kit; no novelty near test windows) &
Week 1 &
Coach Lead &
\textbf{HOLD} gate-grade tests until tool identity and environment unit logging are stable; reslot any compromised test; treat missing fields as \textbf{INVALID} for gate use. \\

\textbf{DEP-1C-013} &
7.07 Load-Bearing Rucking and Carry Systems — Tier 2 (verified dry-load capability; stable ballast; strap reliability) &
Week 4 &
Equipment Lead &
Do not attempt \textbf{MET-2B-021}; reslot long ruck milestone; execute only supporting low-risk exposure until ruck system is verified and stable. \\

\textbf{DEP-1C-007} \& \textbf{DEP-1C-008} &
7.11 Footwear Systems and Foot Care — Tier 3 (rotation pairs; drying tools; blister protocol enforced) &
Week 4 &
Equipment Lead &
\textbf{HOLD} long ruck milestone; downshift to lower-friction modalities; treat any hotspot escalation as a durability constraint; reslot gait-dependent tests. \\

\textbf{DEP-1C-022} &
7.19 Protective and Assistive Equipment — Tier 3 (enhanced blister/chafe interface protection capability) &
Week 5 &
Equipment Lead &
Reduce interface risk: downshift exposure; prohibit “prove it” long ruck attempts; reslot \textbf{MET-2B-021} until protective capability is established. \\

\textbf{DEP-1C-019} &
7.14 Environmental Apparel, Cold, Hot, Wet, Inclement Weather — Tier 3 (seasonal extremes readiness) &
Week 5 &
Trainee &
If environment hazards cannot be controlled safely, substitute indoors where comparable; otherwise reslot tests; do not attempt outdoor gate tests under hazard conditions. \\

\textbf{DEP-1C-011} &
7.17 Aquatic Training and Water Safety Equipment — Tier 2 (pool validity and safety prerequisites for surface swim only) &
Week 10 &
Coach Lead &
If pool validity cannot be verified, do not attempt swim test; reslot \textbf{SWIM} gate item to next protected window; no underwater substitution. \\

\end{longtblr}

\endgroup

\subsection*{14. Validity Requirements}
\phantomsection


\begin{itemize}
\item Evidence admissibility is governed by Stage 1.F and Stage 2 validity rules; only \textbf{VALID} test evidence supports gate outcomes.
\item Missing required Stage 2.E fields (environment/tool identity, warm-up logging when required, measurement context) defaults a test to \textbf{INVALID} for gate use.
\item \textbf{INVALID} tests provide no gate credit and require reslot to the next protected window per Stage 3.E; do not compress by stacking make-up events.
\item \textbf{CAL} events require compliant artifact evidence for \textbf{VALID} gate use; missing/non-compliant artifacts yield \textbf{INVALID} for gate use even if reps/time were achieved.
\end{itemize}

\end{landscape}

\subsection*{15. Change Control Notes}
\phantomsection


\begin{itemize}
\item Allowed without patch:
  \begin{itemize}
  \item within-window sequencing adjustments that preserve test constructs and preserve spacing intent,
  \item reslotting a missed/invalid test to the next protected Deload+Test window within the phase, per Stage 3.E.
  \end{itemize}
\item Requires patch (Stage 1.F change control):
  \begin{itemize}
  \item adding a new test week,
  \item extending phase duration,
  \item changing gate test set or substituting instruments,
  \item changing work-to-deload pattern.
  \end{itemize}
\item Freeze posture near protected windows:
  \begin{itemize}
  \item no novelty in footwear, route class, ruck configuration, pool unit assumptions, or timing tools inside protected windows; violations render evidence non-admissible.
  \end{itemize}
\end{itemize}

\subsection*{16. Compliance Checklist}
\phantomsection


\begin{itemize}
\item Tests listed are Stage 2-defined (C\# and \textbf{MET-2B}-\#\#\#) and limited to major gate-critical items.
\item Deload/Test weeks are identified and aligned to the taper-aware Phase 08 placement posture, with Stage 3 controlling final placement.
\item Dependencies are declared using 7.xx category IDs, due-by week-in-phase, owners, and operational fail actions; 7.02 is not required.
\item Weekly structure table is complete and Cardio cell structure includes Modality; Minutes; RPE plus Talk Test cue; Continuity continuous planned.
\item Operating constraints state: 6 sessions/week, required elements order, Cardio minimum standard in body text, and 10,000 steps/day in body text.
\item Validity posture is explicit: \textbf{VALID} evidence only for gates; \textbf{INVALID} yields no gate credit and triggers reslot per Stage 3.E.
\item Governing Status is justified as Deployable: no placeholders appear in Mid-Phase Gate, End-Phase Gate, Test Battery, Exit Standards, or Dependencies affecting gating.
\end{itemize}

\end{document}